\documentclass[11pt]{article}
\usepackage[margin=1.1in]{geometry}
\usepackage[cmex10]{amsmath}
\usepackage{amssymb,amsthm}
\usepackage[T1]{fontenc}
\usepackage[hidelinks,hypertexnames=false]{hyperref}
\allowdisplaybreaks[2]
\setlength{\emergencystretch}{2em}
\title{Simple negative zeros for very tall balanced arrays}
\author{ProgramComputer}
\date{}
\newtheorem{theorem}{Theorem}
\DeclareFontShape{OMX}{cmex}{m}{n}{<-> cmex10}{}
\begin{document}
\maketitle

\begin{abstract}
Let $P_n^{(r,d)}$ enumerate labeled binary $n\times d$ arrays with row sums
at most $r$ by their common column sum. The theorem gives a single existential
constant $C$ for which every $P_n^{(r,d)}$ with $d\ge2$, $1\le r\le d$
and $n\ge Cd^4\log(d+1)$ has only simple, strictly negative zeros.
The proof combines adjacent coefficient estimates near the two ends
with an evaluation of the balance integral at a pair of dominant phases.
\end{abstract}

\section{Statement}
The \href{https://www.symmetricfunctions.com/realRootedOpenProblems.htm\#conjBalancedArrays}{balanced-array conjecture of Per Alexandersson}
asks for real-rootedness throughout the parameter range.
For the labeled arrays in the abstract let $a_k$ count those with
each column sum equal to $k$, and write
\[
P_n^{(r,d)}(q)=\sum_{k=0}^{\lfloor nr/d\rfloor}a_kq^k.
\]
\begin{theorem}
There exists $C>0$ such that for every integer $d\ge2$, every $1\le r\le d$,
and every integer $n\ge Cd^4\log(d+1)$, the polynomial $P_n^{(r,d)}$
has $\lfloor nr/d\rfloor$ simple, strictly negative zeros.
\end{theorem}

Thus a counterexample to the all-parameter real-rootedness conjecture must have $n<Cd^4\log(d+1)$ for the existential absolute constant $C$ in the theorem. The original conjecture remains open.

\section{Reduction and angular levels}

Failure for every $C$ would give a sequence with $n/(d^4\log(d+1))\to\infty$. Bounded $d$ is excluded by the fixed-cap theorem, as finitely many caps are then involved. The full cap $r=d\ge2$ has simple negative roots by the binomial multiplier theorem. Fix the absolute near-saturation threshold $h_0$ from the upper saddle estimate first. The interior estimate then has fixed positive deficit parameter $\eta=h_0$, and its cutoff constant can be chosen after $h_0$ is fixed. Pass to $d\to\infty$ and $r<d$. Put $T=nr/d$, $D=\lfloor T\rfloor$, $b=d^2/n$ and $\nu=(d-1)/2$. We have $b\to0$ and $D/d\to\infty$. Every integer quota between 0 and $D$ is feasible. For instance, split $dk$ slots among the $n$ rows in groups of at most $r$ and fill the slots successively with the repeating sequence of columns $1,\ldots,d$. \href{https://github.com/ProgramComputer/open-math-research/blob/f7e10f84dd21fec44ca58c60bae0879d12ebd89b/topics/balanced-arrays/finite-linear-reduction.md}{Full-cap result}

At row tilt $w>0$ put $H(w)=\sum_{s=0}^r\binom ds w^s$, with normalized size law $S$, mean $\mu$, variance $V$, and $h=r-\mu$. Write $w_x$ for the tilt with $y(w_x)=x$, where $y=n\mu/d$. This covers $0<x<T$ since $\partial_{\log w}\mu=V>0$. For $t_*=\pi/d$ and $G(s)=s(d-s)/(d(d-1))$, define

\begin{equation*}
\lambda_t=\frac{\mathbb E[G(S)e^{itS}]}{\mathbb E e^{itS}},
\qquad \Psi(w)=n\arg H(we^{it_*})-\nu\arg\lambda_{t_*}.
\end{equation*}

At $0\le t\le t_*$ both positive-weight characteristic sums lie in the closed sector from zero angle to $rt<\pi$, so the arguments continued from zero exist. The limiting values of $\Psi$ at zero and infinity are $-\nu t_*$ and $T\pi$. The lower limit uses the first nonzero degree of the $G$-weighted polynomial, and the upper limit uses $G(r)>0$. For $1\le k<D$ define $v_k$ as the first radius attaining $\Psi=k\pi$. Continuity gives strictly increasing $v_k$.

The following estimate controls the locations of the levels with an unbounded cap:

\begin{equation}
|\Psi(w)/\pi-y|\le C_1 d+\frac{C_1}{\sqrt d}\min(y,T-y).
\tag{M1}
\end{equation}

For a coordinate indicator $X_1$ and $p=\mu/d$, exchangeability gives $V=d\,\mathrm{Cov}(S,X_1)=dp(1-p)(\mathbb E[S\mid X_1=1]-\mathbb E[S\mid X_1=0])$. Given $X_1=1$ or $X_1=0$, the number of other successes has the same $(d-1)$-trial tilted binomial law truncated at $r-1$ or $r$, respectively. The former truncation has smaller mean, so the conditional difference for $S$ is at most one and $V\le\mu(1-\mu/d)\le d/4$. Also $V\le m(1+m)$ for $m=\mu$ and for $m=h$. To see this for any nonnegative integer variable $X$ with log-concave mass on $[0,r]$ and mean $m$, write $F_j=\mathbb P(X\ge j)$. Its ratios $F_{j+1}/F_j$ decrease since $F_j/\mathbb P(X=j)$ is the sum of products of consecutive decreasing mass ratios. Hence $F_{i+j}\le F_iF_j$ and

\begin{equation*}
\mathbb E[X(X-1)]/2=\sum_{i,j\ge1}F_{i+j}\le m^2 .
\end{equation*}

Apply this to $S$ and $r-S$. Combining the two variance bounds at $m=\sqrt d$ gives $V\le C\sqrt d\min(\mu,h)$. Now $|\mathbb E e^{it(S-\mu)}-1|\le V t^2/2$ by Taylor with centered linear term. This is $O(1/d)$ on $[0,t_*]$, whence $|\arg H(we^{it_*})-\mu t_*|\le C V/d^2$. Finally $|\arg\lambda_{t_*}|\le\pi$ by the sector property. Substituting gives (M1). \hyperref[sec:interior-saddle]{Uniform interior saddle and variance bound}

Consequently, at $w=v_k$, bounded $k/d$ implies $y=O(d)$, and bounded $(D-k)/d$ implies $T-y=O(d)$. If both indices divided by $d$ tend to infinity, so do $y/d$ and $(T-y)/d$, by applying (M1) to either hypothetical bounded subsequence.

\section{The determinant phase at the ends}

Suppose $k/d\to c$ with $0<c<\infty$. By (M1), $\mu=O(b)$ at the crossing. The size law then has $\mathbb E|S-\mu|^3=O(\mu)$, and the law reweighted by $G$ has expected distance $O(\mu)$ from size one. Indeed $dw=O(\mu)$, $p_s\le(dw)^s/s!$, $p_1\asymp\mu$, and $G(s)/G(1)\le s$. Thus $\arg\lambda_{t_*}=t_*+O(\mu t_*)$. The centered characteristic argument at $t_*$ is $O(\mu t_*^3)$ by the vanishing linear term and the third absolute moment. Since $n\mu t_*^3=O(y/d^2)=o(1)$, the level equation gives $y(v_k)=k+1/2+o(1)$.

Suppose instead $(D-k)/d\to c$ with $0<c<\infty$. We have $h=O(b)$ at the crossing. For $J=r-S$ let $p'_j=\mathbb P(J=j)$. By log-concavity all forward ratios are at most $p'_1/p'_0\le h/(1-h)$. It follows that $\mathbb E|S-\mu|^3=O(h)$, and that the $G$-weighted mean distance from $r$ is $O(h)$, using $G(r-j)/G(r)\le1+j$. We obtain $\arg\lambda_{t_*}=O(h t_*)$ and a centered characteristic argument $O(h t_*^3)$. Thus $y(v_k)=k+o(1)$.

At real $x$ within bounded distance of either such $k$, the variance is asymptotic to $dx/n$ at the lower end and to $d(T-x)/n$ at the upper end. At the lower end this follows from the small-mean factorial moments. At the upper end it follows from the geometric domination of $J$, which gives $\mathbb E J(J-1)=O(h^2)$. Differentiation yields $\frac{\mathrm d}{\mathrm dx}(d\log w_x)=b/V_x$, tending to $1/c$ uniformly on each bounded window. Integrating from $k$ to $y(v_k)$ gives the two expansions

\begin{equation}
d\log v_k=d\log w_k+d/(2k)+o(1)\quad(k/d\to c),
\qquad d\log v_k=d\log w_k+o(1)\quad((D-k)/d\to c).
\tag{M2}
\end{equation}

\section{Signs and degree}

Write $t_k=\sqrt{a_{k-1}/a_{k+1}}$ and $\Gamma_k=a_k^2/(a_{k-1}a_{k+1})$. In either finite positive overlap above, endpoint coefficient estimates show that $\log\Gamma_{k+j}\to1/c$ for every fixed integer $j$ and $\log t_k=d\log v_k+o(1)$. The lower endpoint proof compares with independent columns. The upper endpoint proof conditions on the complete row-size profile, including at exact saturation. At distances with $k/d\to0$ or $(D-k)/d\to0$, they show $\Gamma_k\to\infty$. \hyperref[sec:nearly-empty]{Nearly empty rows} \enspace \hyperref[sec:upper-quota]{Nearly saturated coefficients}

By the orientation inequality the $a_k$ are strictly log-concave at internal quotas. If $\Gamma_k>4$, the centered normalized alternating polynomial at $t_k$ is at least $1-2/\sqrt{\Gamma_k}>0$ because the two tails beginning at distance two alternate with decreasing magnitudes. In the finite positive overlaps, for $x=t_k$ or $v_k^d$, the terms at each fixed distance $j$ from the center converge after normalization to $e^{-j^2/(2c)}$. Indeed the terms at distance one have the required relative magnitudes by the definition of $t_k$ and (M2), and the consecutive logarithmic differences follow from $\log\Gamma_{k+j}$. Beyond any large fixed distance $L$, log-concavity bounds the outward ratios by a geometric factor $<1$ uniformly, using the ratio at $L$. Taking $L\to\infty$ therefore yields

\begin{equation*}
\frac{(-1)^kP(-x)}{a_kx^k}\longrightarrow
\sum_{j\in\mathbb Z}(-1)^j e^{-j^2/(2c)}
=\sqrt{2\pi c}\sum_{m\in\mathbb Z}e^{-2\pi^2 c(m+1/2)^2}>0.
\end{equation*}

Here the last identity is Poisson summation. \href{https://github.com/ProgramComputer/open-math-research/blob/f7e10f84dd21fec44ca58c60bae0879d12ebd89b/topics/balanced-arrays/orientation-normalization.md}{Orientation inequality}

If $k/d,(D-k)/d\to\infty$, set $w=v_k$. Then $n\mu/d^2,nh/d^2\to\infty$ from (M1). On the product-negative torus, character integration of the $n$th power of the row generating function gives $P(-w^d)$, as only the common column exponents survive. The two principal contributions have phase $\pm\Psi(w)$ and equal positive amplitude. Their sum differs from the exact polynomial by $o(1)$ times the amplitude of one contribution. This conclusion covers every tilt: use the small-mean proof for $\mu\le1/10$, the near-full proof for $h\le h_0$ with fixed absolute $h_0>0$, and the interior proof for $\mu\ge1/10$, $h\ge h_0$. All three estimates use the branch of the quadratic determinant obtained from the positive real tilt. At the first crossing $\cos\Psi(w)=(-1)^k$, giving the desired strict sign. \hyperref[sec:nearly-empty]{Nearly empty rows} \enspace \hyperref[sec:growing-cap-upper-saddle]{Near-full angular estimate} \enspace \hyperref[sec:interior-saddle]{Interior angular estimate}

For $k\le d$ or $D-k\le d$ use $t_k$, and at all other internal quotas use $v_k^d$. Inside their respective ranges the points increase by strict log-concavity and by continuity at the first crossings. At either join the adjacent geometric logarithms differ by $(\log\Gamma_k+\log\Gamma_{k+1})/2\to1$; replacing one geometric radius by the corresponding phase radius changes its logarithm by $o(1)$ by the overlap estimates. The endpoints are zero and positive infinity, with signs $1$ and $(-1)^D$. A subsequence with any failed interior sign has a further subsequence with limits of $k/d$ and $(D-k)/d$ in $[0,\infty]$. Their sum tends to infinity. The arguments above exclude a zero limit, a positive finite limit and two infinite limits. Thus the signs alternate at every one of these ordered points. By continuity the reflected polynomial has $D$ distinct positive zeros, exhausting its degree. This gives the stated simplicity and negativity and rules out the assumed sequence of counterexamples.


\appendix
\clearpage
\section{Upper quotas with a growing row cap in very tall arrays}\label{sec:upper-quota}

Let $d\to\infty$, $n\ge C_*d^4\log(d+1)$ for a sufficiently large absolute constant $C_*$ and $1\le r<d$, allowing $r$ to vary. Write $T=nr/d$, $D=\lfloor T\rfloor$, $H(w)=\sum_{s=0}^r\binom ds w^s$, and $A_i=[w^{di}]H(w)^n$. Let $a_i$ count binary arrays with all column sums $i$ and all row sums at most $r$. Uniformly for $0\le D-i\le Cd$ at any fixed $C$, we have

\begin{equation*}
a_i=A_i\sqrt d\,(2\pi\Lambda_*(i))^{-(d-1)/2}(1+o(1)),
\qquad
\Lambda_*(i)=\frac{i(d-r)}{d-1}.
\end{equation*}

At internal quotas set $\Gamma_k=a_k^2/(a_{k-1}a_{k+1})$ and $t_k=(a_{k-1}/a_{k+1})^{1/2}$. If $(D-k)/d\to c\in(0,\infty)$, then $\log\Gamma_{k+j}\to1/c$ for any fixed integer $j$ and $\log t_k=d\log w_k+o(1)$, where $w_k$ gives the size law of mean $dk/n$. If $(D-k)/d\to0$, then $\Gamma_k\to\infty$.

\subsection{Fourier balance uniform over upper profiles}

Fix a profile $(s_1,\ldots,s_n)$ with $0\le s_j\le r$ and total $di$, and set $e=nr-di=O(d^2)$. At least $n-e$ rows have size $r$. For each size $s$, let $f_s(\theta)=\mathbb E e^{i\sum_{a\in U}\theta_a}$ for a uniformly chosen $s$-subset $U$ of the columns. The balance probability is the pinned $(d-1)$-torus integral of $\prod_jf_{s_j}(\theta)e^{-\mathrm{i}\,i\sum_a\theta_a}$, with $\theta_d=0$ and normalized Haar measure. Common shifts leave the integrand unchanged because the prescribed total is $di$.

Set $V=\mathbf1^\perp\subset\mathbb R^d$ and $R=d^{-1}$. Project onto $V$ along the common angle. The Euclidean ball $\|y\|\le R$ in $V$ is injective on the resulting quotient: its pinned inverse is $\theta_a=y_a-y_d$ modulo $2\pi$, with all lifts of magnitude at most $2R$. The Jacobian from pinned coordinates to Euclidean hyperplane volume is $1/\sqrt d$, since the projection of their basis has Gram matrix $I_{d-1}-\mathbf1\mathbf1^\top/d$. Normalized Haar measure on the ball is therefore $\sqrt d\,(2\pi)^{-(d-1)}dy$.

Put $z_j=e^{\mathrm{i}\theta_j}$ and $D_z=\sum_j|z_j-\bar z|^2$. Generate a uniform $s$-subset by placing a uniform permutation into slots, tossing a fair coin to select one member from each of the first $m=\min(s,d-s)$ disjoint pairs, and taking $s-m$ fixed unused slots. Conditional Fourier modulus is at most $e^{-X/8}$, where $X$ sums the squared chords of the pairs, since each fair choice has modulus $\sqrt{1-|z_i-z_j|^2/4}$. A random distinct pair has mean square chord $2D_z/(d-1)$, hence $\mathbb E X=2mD_z/(d-1)$. Each chord square is at most four. For two disjoint ordered pairs the expected product is at most $d(d-1)/[(d-2)(d-3)]$ times the square of the one-pair mean, by allowing intersections in the numerator. Thus $\mathbb E X^2\le C(\mathbb E X+(\mathbb E X)^2)$. The inequality $1-e^{-X/8}\ge cX/(1+X)$ and Cauchy--Schwarz give a modulus loss at least $c\min(mD_z/d,1)$, hence at least $ca_s\min(D_z,1)$ with $a_s=s(d-s)/[d(d-1)]$.

Outside the projected ball $D_z\ge c_0R^2$. Indeed anchor at an index minimizing $|z_j-\bar z|$. The sum of squared chords to that anchor is at most $2D_z$, so the phases relative to it in $[-\pi,\pi]$ have squared Euclidean norm at most $C D_z$ by the chord-angle inequality. Removing their average can only decrease that norm. If $D_z<c_0R^2$ for small enough $c_0$, the centered representative lies inside the ball. Applying the modulus bound to the $n-e$ size-$r$ rows, with $a_r\ge1/d$, bounds the exterior integral by $\exp(-c'n/d^3)$.

On the ball let $L=\sum_{a\in U}y_a$ for a row of size $s$. By exchangeability $\mathbb E L=0$ and $\mathbb E L^2=a_s\|y\|^2$. For the fourth moment replace the subset by its complement if $s>d/2$, which negates $L$. Set $p_j=(m)_j/(d)_j$ and $Q_j=\sum_a y_a^j$. Expansion using $Q_1=0$ gives

\begin{equation*}
\mathbb E L^4=(p_1-7p_2+12p_3-6p_4)Q_4+(3p_2-6p_3+3p_4)Q_2^2.
\end{equation*}

The index patterns $4,3+1,2+2,2+1+1$ give $p_1Q_4$, $-4p_2Q_4$, $3p_2(Q_2^2-Q_4)$ and $6p_3(2Q_4-Q_2^2)$. Subtracting these from $(\sum_a y_a)^4=0$ with all $p_j$ set to one gives the distinct-index term $p_4(3Q_2^2-6Q_4)$. Since $p_j\le p_1\le2a_s$, the fourth moment is at most $Ca_s\|y\|^4$. Cauchy--Schwarz yields $\mathbb E|L|^3\le Ca_s\|y\|^3$. Taylor's inequality and the logarithm series therefore give

\begin{equation*}
\log f_s(y)=-\frac{a_s}{2}\|y\|^2+O(a_s\|y\|^3).
\end{equation*}

The branch passes through zero at $y=0$. The extra series error $O(a_s^2\|y\|^4)$ is absorbed since $a_s\le1$ and $\|y\|\le R$. Thus the centered integrand is $\exp(-\Lambda\|y\|^2/2+B(y))$, with

\begin{equation*}
\Lambda=\sum_j a_{s_j},\qquad |B(y)|\le C\Lambda\|y\|^3,\qquad
(n-e)/d\le\Lambda\le n.
\end{equation*}

Use $|e^B-1|\le |B|e^{|B|}$ and $\|y\|\le1/d$. The relative error integrated over the ball, normalized by the full target Gaussian integral, is at most

\begin{equation*}
C\Lambda(1-2C/d)^{-(d-1)/2}\mathbb E\|Y\|^3
=O(d^{3/2}/\sqrt\Lambda)=O(d^2/\sqrt n)=o(1),
\end{equation*}

where $Y$ is Gaussian on $V$ of precision $\Lambda(1-2C/d)$. The normalization ratio is bounded and the moment bound follows from the Gaussian fourth moment. The target Gaussian mass outside the ball is $O(d/(\Lambda R^2))=O(d^4/n)=o(1)$ by Markov's inequality. Its full torus-density-normalized integral is $\sqrt d(2\pi\Lambda)^{-(d-1)/2}\ge e^{-C_1d\log n}$ for fixed $C_1$. As $n/\log n$ increases for $n>e$, the ratio $n/(d^4\log n)$ has lower limit at least $C_*/4$ in the stated range. Choose $C_*$ large enough that $c'n/d^3\ge 2C_1 d\log n$ there for all large $d$. The exterior is then negligible relative to the Gaussian. All estimates are uniform in the profile and in $r$.

The profile dependence in the determinant can now be discarded, since

\begin{equation*}
0\le \Lambda-\Lambda_*(i)
=\frac{\sum_j s_j(r-s_j)}{d(d-1)}\le\frac{re}{d(d-1)},
\qquad
d\log(\Lambda/\Lambda_*(i))
\le\frac{dre}{di(d-r)}
=O(d^3/n).
\end{equation*}

In the last estimate use $di=nr-e\sim nr$ and $d-r\ge1$. Averaging over profiles in $A_i$ with weights proportional to $\prod_j\binom d{s_j}$ gives the asserted count, including $e=0$. The profile event describes exactly the balanced arrays conditional on the row sizes. The logarithmic first differences of the determinant factor at these quotas are $O(d/i)=O(d^2/(nr))=o(1)$, so the same is true of its second differences. Together with the uniform relative error, this means all required fixed-offset logarithmic differences may be taken on $A_i$.

\subsection{Scalar quota differences}

Put $b=d^2/n$. At a real $x\in(0,T)$, let the size law with positive tilt $w_x$ have mean $dx/n$. The mean is strictly increasing in $\log w_x$ since its derivative is the positive variance $V_x$. As $w_x$ spans the positive axis the mean spans $(0,r)$. The deficit $J=r-S$ has mean $h_x=d(T-x)/n$ and weights $B_j z^j$ for $B_j=\binom d{r-j}$ and $z=1/w_x$. Write $\ell=B_1z/B_0$. The forward weight ratios decrease in $j$, so $B_j z^j/B_0\le\ell^j$. If $h_x=o(1)$, the mass at zero is at least $1-h_x$ while the mass at one is at most $h_x$, so $\ell\le h_x/(1-h_x)$. Summing the resulting geometric majorant and retaining the first two masses gives, uniformly in the cap,

\begin{equation*}
h_x=\ell+O(\ell^2),\qquad
V_x=h_x+O(h_x^2),\qquad
\mathbb E|J-h_x|^3=O(h_x),\qquad
\mathbb P(J=0)\mathbb P(J=1)\asymp h_x.
\end{equation*}

Here all raw moments up to order three from $J\ge2$ are $O(\ell^2)$, since $\sum_{j\ge2}j^m\ell^j=O(\ell^2)$ for $\ell\le1/2$. The derivative of $V_x$ with respect to the size mean is the signed third centered size moment divided by variance, hence is uniformly $O(1)$ in this regime.

Suppose $(D-k)/d\to c\in(0,\infty)$. On any bounded real window around $k$, $h_x\sim cb$ and $nh_x\sim cd^2$. At integer $i$ in that window, the sum of $n$ tilted deficits has integral mean $e=nr-di$. The two-mass product estimate implies that the modulus of the deficit characteristic function is at most $\exp(-c_2h_i t^2)$ for $|t|\le\pi$ by comparing two independent copies and retaining their differences $\pm1$. Rescale Fourier inversion at the mean by $t=v/\sqrt{nV_i}$. The centered expansion of the logarithm on bounded $v$ gives $-v^2/2+O(|v|^3/\sqrt{nh_i}+v^4/n)$ after multiplication by $n$. The modulus bound supplies an integrable Gaussian majorant independent of the parameters. Thus the central probability is $(2\pi nV_i)^{-1/2}(1+o(1))$ at each fixed integer shift.

For $E(x)=n\log H(w_x)-dx\log w_x$, one has $E'=-d\log w_x$ and $E''=-b/V_x\to-1/c$ uniformly on the window. Also $E'''=O(b(d/n)/h_x^2)=O(1/d)$ by the derivative bound just proved. The tilt identity yields $\log A_i=E(i)-\tfrac12\log(2\pi nV_i)+o(1)$, and neighboring $V_i$ have ratio tending to one. The central second difference of $E$ gives $\log\Gamma_{k+j}\to1/c$, and $(E(k-1)-E(k+1))/2=-E'(k)+O(1/d)$ gives the claimed logarithm of $t_k$. The determinant factor has already been accounted for.

\subsection{Distances sublinear in dimension}

Use scalar deficit counts $U_e=[z^e](\sum_j B_j z^j)^n$. They are log-concave relative to $\binom{n+e-1}e$. To see this, let $C_m$ be the analogous totals for $n-1$ rows, extended by zero. They are log-concave with interval support by convolution. For increasing $e$, the law on $j$ with weights $(j+1)B_j C_{e-j}$ increases in likelihood ratio order: the common-support ratio is proportional to $C_{e+1-j}/C_{e-j}$, increasing with $j$, and the support endpoints both increase weakly. Under this law the expectation of the decreasing sequence $B_{j+1}/B_j$, extended by zero at $r$, equals $(e+1)U_{e+1}/[(n+e)U_e]$. The numerator comes from differentiating the row polynomial power, and the denominator from exchangeability of the rows conditional on their sum. Thus the normalized forward ratio decreases, proving the assertion.

At an internal quota with $(D-k)/d\to0$, $e=d(T-k)\ge d$ and $e=o(d^2)$. The log curvature of $A_k=U_e$ at quota step one is at least the log curvature of its negative-binomial normalizer at deficit step $d$. Writing the binomial forward logarithmic increment as $g(v)=\log[(n+v-1)/v]$, this gives

\begin{equation*}
\log\frac{A_k^2}{A_{k-1}A_{k+1}}
\ge\sum_{j=1}^d \big(g(e-d+j)-g(e+j)\big)
\ge\frac{d^2(n-1)}{(e+d)(n+e+d)}\longrightarrow\infty .
\end{equation*}

Indeed $-g'(v)=(n-1)/[v(n+v-1)]$ and each difference integrates over length $d$. This argument includes the case $e=d$, where its upper-quotient neighbor has exact saturation. The determinant comparison transfers the divergence to $\log\Gamma_k$.

\section{Nearly empty rows, uniformly in the cap}\label{sec:nearly-empty}

Let $1\le r<d$, $d\to\infty$, and let the row-size law be proportional to $\binom ds w^s$ on $0\le s\le r$. Write $\mu=\mathbb E S$, $V=\mathrm{Var}(S)$ and $H(w)=\sum_{s\le r}\binom ds w^s$. The following two estimates use constants independent of $r$.

\subsection{Two principal phases at small mean}

If $0<\mu\le1/10$ and $n\mu/d^2\to\infty$, the torus integral of the normalized row function to the power $n$ at the negative argument $-w^d$ has the same two-saddle formula as in the interior estimate. In detail, write $\phi=\pi/d$, $G(s)=s(d-s)/(d(d-1))$, $\lambda_t=\mathbb E[G(S)e^{itS}]/\mathbb E e^{itS}$, and $\nu=(d-1)/2$. Its value is

\begin{equation*}
\frac{2}{\sqrt d}(2\pi n)^{-\nu}
\mathrm{Re}\!\left[(\mathbb E e^{i\phi S})^n\lambda_\phi^{-\nu}\right]
+o\!\left(\frac{|\mathbb E e^{i\phi S}|^n}{\sqrt d}(2\pi n|\lambda_\phi|)^{-\nu}\right).
\end{equation*}

We detail the uniformity needed for the moving-tilt and transverse integration lemmas. Write $a=dw=p_1/p_0$ where $p_j=\mathbb P(S=j)$. Since $p_0\ge1-\mu$ and $p_1\le\mu$, $a\le1/9$. The bounds $p_s\le a^s/s!$ and $p_1\ge(9/10)a$ imply $a\asymp\mu$, $p_1\asymp\mu$, $V\asymp\mu$ and $\mathbb E S^m\le C_m\mu$ for each fixed $m\ge1$. Also $V=\mu+O(\mu^2)$ as $\mu\to0$. For each $t$, $|\mathbb E e^{itS}|\ge1-2\mu\ge4/5$. The empty support and each singleton give the angular bound

\begin{equation*}
|R(\theta)/H(w)|\le\exp\!\left[-c\,\frac{\mu}{d}\sum_{j=1}^d|e^{i\theta_j}-1|^2\right],
\qquad R(\theta)=\sum_{|A|\le r}w^{|A|} e^{i\sum_{j\in A}\theta_j}.
\end{equation*}

Indeed these terms in the squared characteristic modulus give a loss $p_0p_1 d^{-1}\sum_j|e^{i\theta_j}-1|^2$. Outside a fixed small arc around zero in any coordinate, retain half its exponential loss and integrate the remaining $d-1$ free circle variables (eliminating a different coordinate). The absolute bound, divided by $H(w)^n$, is $d e^{-c n\mu/d}(C d/(n\mu))^{(d-1)/2}$. At the principal phase $|\mathbb E e^{i\phi S}|\ge e^{-C\mu/d^2}$ and $\lambda_0\asymp\mu/d$ using singleton mass and $G(s)\le s/d$. Thus the off-arc bound is negligible relative to one saddle whenever $n\mu/d^2\to\infty$.

Inside the arc, the first $d-1$ independent angles determine the odd common phase $t=(2j+1)\pi/d$ by averaging all lifted coordinates. Write the zero-sum deviations as $y$. The moment bounds used in the fixed-cap local and global integral estimates hold with constants independent of $r$. Namely, for $L=\sum_{j\in I} y_j$ at a uniform size-$S$ support $I$,

\begin{equation*}
\mathbb E[S^l|L|^m]\le \frac{\mathbb E S^{m+l}}{d}\sum_j|y_j|^m
\le C_{m,l}\frac{\mu}{d}\sum_j |y_j|^m
\qquad(m\ge1,\ l\ge0).
\end{equation*}

The first inequality conditions on $S$ and uses convexity. The same inequalities bound any fixed number of derivatives in $t$ of the complex-tilted moments, since the denominator stays above $4/5$. In particular $\mathrm{Re}\lambda_t=\lambda_0(1+O(t^2))$, the cubic term is $-\mathrm{i}c(t)\sum_j y_j^3/6$, where $|c(t)|=O(\mu/d)$ and $|\mathrm{Im}\,c(t)|=O(\mu|t|/d)$. Its real part is therefore $O((\mu/d)|t|\sum_j|y_j|^3)$, and the log expansion's quartic remainder is $O((\mu/d)\sum_j y_j^4)$. The central characteristic expansion gives $(\log|\mathbb E e^{itS}|)''=-V+O(\mu |t|)\le-c\mu$ in a sufficiently small fixed arc, by the centered moments just established. These are precisely the bounds used for the one-dimensional convolution majorant and the complex Gaussian approximation on the centered hyperplane. The convolution bound makes the sum of nonprincipal clusters relative to one saddle at most

\begin{equation*}
C\exp(C d^2/(n\mu))\sum_{j\ge1}\exp[-c(n\mu/d^2)((2j+1)^2-1)]=o(1).
\end{equation*}

The relative error on either principal cube is $O(d/\sqrt{n\mu}+d^2/(n\mu))$ by the integrated Wick estimate for the centered Gaussian; there $d\phi^2=O(1/d)$. Choose the singleton zero arc $|\theta_j|\le\eta$ inside the small-phase range of the convolution estimate, reducing the fixed $\eta$ if necessary. Any lifted cluster has common phase in $[-\eta,\eta]$ and centered coordinates of magnitude at most $2\eta$. At either principal phase $\pm\pi/d$, the fixed centered cube of radius $\eta/2$ lies within the original zero arc for all sufficiently large $d$. The corresponding principal region therefore lies between cubes of radii $\eta/2$ and $2\eta$ in the centered hyperplane. The flagged-coordinate convolution estimate controls its part outside the smaller cube. Thus all constants in the two cited estimates are uniform here even for an unbounded cap. \href{https://github.com/ProgramComputer/open-math-research/blob/f7e10f84dd21fec44ca58c60bae0879d12ebd89b/topics/balanced-arrays/fixed-cap/fixed-cap-moving-u-localization.md}{Convolution estimate and cluster coordinates} \enspace \href{https://github.com/ProgramComputer/open-math-research/blob/f7e10f84dd21fec44ca58c60bae0879d12ebd89b/topics/balanced-arrays/fixed-cap/tall-sparse-transverse-saddle.md}{Integrated local Gaussian estimate}

\subsection{Coefficients at positive finite scaled quotas}

Suppose now $n/d^4\to\infty$. In the unrestricted column model, choose independently for each column a uniform $i$-subset of the $n$ rows. The expected number of unordered pairs of occupied cells sharing a row is $\binom d2 i^2/n$. For $i=k+O(1)$ this is $O(d^4/n)=o(1)$. If no pair occurs, the cap is satisfied for every $r\ge1$. Therefore

\begin{equation*}
a_i=\binom ni^d(1+o(1))
\end{equation*}

uniformly on every fixed window. Thus, with $\Gamma_i=a_i^2/(a_{i-1}a_{i+1})$ and $t_i=\sqrt{a_{i-1}/a_{i+1}}$, we have $\log\Gamma_{k+j}\to1/c$ for each fixed integer $j$. Also

\begin{equation*}
\log t_k=d\log\frac{k}{n-k}+\frac d2\log(1+1/k)-\frac d2\log(1+1/(n-k))+o(1).
\end{equation*}

At the scalar tilt $w_k$ prescribed by $\mathbb E S=dk/n$, the small-mean series gives $dw_k=\mu+O(\mu^2)$. Hence $d\log(w_k/(k/(n-k)))=O(d\mu+dk/n)=o(1)$. We obtain $\log t_k=d\log w_k+d/(2k)+o(1)$ for these quotas. For $k/d\to0$ the orientation inequality gives $\Gamma_k\ge(1+1/k)^d\to\infty$ directly.

\section{Uniform interior saddle for very tall capped arrays}\label{sec:interior-saddle}

Let $d\to\infty$, $1\le r<d$, $v>0$, and let $S$ have probabilities proportional to $\binom ds v^s$ on $0\le s\le r$. Write $\mu=\mathbb ES$, $h=r-\mu$, $V=\mathrm{Var}(S)$ and $H=\sum_{s\le r}\binom ds v^s$. Assume $\mu\ge 1/10$ and $n\ge C_* d^6\log(d+1)$ with $h\ge d^{-3}$, for a sufficiently large absolute constant $C_*$. Alternatively, fix any $\eta>0$ and assume $h\ge\eta$ and $n\ge C_\eta d^4\log(d+1)$, for a sufficiently large $C_\eta$. Given $S$, choose a uniform $S$-subset $I$. Put $R(\theta)=\mathbb E\exp(i\sum_{i\in I}\theta_i)$, $\phi(t)=\mathbb E e^{itS}$, $a_s=s(d-s)/(d(d-1))$ and

\begin{equation*}
a(t)=\frac{\mathbb E[a_S e^{itS}]}{\phi(t)}.
\end{equation*}

Then $\mathcal I=P_n^{(r,d)}(-v^d)/H^n$ is the Haar integral of $R(\theta)^n$ on the translated torus $\sum_j\theta_j=\pi\pmod{2\pi}$. With $t_*=\pi/d$,

\begin{equation*}
\mathcal I=\frac{2}{\sqrt d}\,(2\pi n)^{-(d-1)/2}
   \mathrm{Re}\!\left[\phi(t_*)^n a(t_*)^{-(d-1)/2}\right]
   +o\!\left(\frac{|\phi(t_*)|^n}{\sqrt d}(2\pi n|a(t_*)|)^{-(d-1)/2}\right).
\end{equation*}

The error is absolute on the scale of one saddle, so the statement remains meaningful if the sum almost cancels. The square root and powers use the branch near positive real $a$. The conclusion is uniform in the cap and tilt within the displayed range. It is a pointwise saddle formula; using it to count roots requires control of the phase as the tilt varies.

\subsection{Tilt estimates and diagonal gap}

Success bias followed by subtracting one gives the $(d-1)$-trial tilt truncated at $r-1$. Failure bias gives the same tilt truncated at $r$. The stochastic comparison of their means gives $V\le\mu(1-\mu/d)\le d/4$. This holds also after any real exponential change of tilt. Integrating the bound for the second derivative of the log moment generating function yields $\mathbb E e^{u(S-\mu)}\le e^{du^2/8}$ for real $u$.

The matching lower bound needed here is $V\ge c\min(1,\mu,h)$. If $V\ge1/100$ this follows immediately. Otherwise the interval of radius $1/3$ about $\mu$ has probability at least $1-9V>0.91$ by Chebyshev and contains at most one integer. Hence an atom at the mode has mass above $0.91$. At an interior mode $m$, the product of adjacent ratios is $p_{m-1}p_{m+1}/p_m^2=m(d-m)/[(m+1)(d-m+1)]\ge 1/4$. Both ratios are at most one, so each is at least $1/4$, contradicting this mass. At mode zero, put $q=p_1/p_0$. Then $q<1/9$, and log-concavity gives $p_j\le p_0 q^j$. Thus $\mu\le p_0 q/(1-q)^2\le Cq$, whereas $V\ge p_0p_1\ge c q$ by the independent-pair variance identity. At mode $r$, apply the same argument to $r-S$. This proves $V\ge c d^{-3}$ in the first regime and $V\ge c\min(1/10,\eta)$ in the second.

Also $\bar a=\mathbb E a_S\ge V/d$ by $dV\le\mu(d-\mu)$. Independently, $a_s\ge1/d$ for every $1\le s\le r<d$, and Cauchy--Schwarz gives $\mathbb P(S\ge1)\ge\mu^2/(\mu^2+V)\ge\mu/(\mu+1)\ge1/11$. Consequently $\bar a\ge1/(11d)$.

To compare diagonal phases, let $\Delta=S-S'$ with independent copy $S'$. Its distribution is symmetric and log-concave, by preservation of discrete log-concavity under convolution. Its masses decrease with $|j|$. Such a law is a mixture of the uniform laws on integer intervals $[-L,L]$: if the mass at $j$ is $g(j)$, the mixing mass at $L$ is $(2L+1)(g(L)-g(L+1))$. Pick $L_0=C\sqrt{d\log d}$ with $C$ an adequately large fixed constant. The sub-Gaussian bound and $|\Delta|\le d$ give $\mathbb E[\Delta^2;|\Delta|>L_0]=o(V)$ uniformly. The part of the mixture at $L>2L_0$ contributes $o(V)$ to the second moment, because under each such uniform law a fixed fraction of its second moment lies above $L_0$ in absolute value.

For the uniform law on $[-L,L]$, $L\ge1$, $1-\mathbb E\cos(t\Delta)\ge c\min(L^2t^2,1)$ for $|t|\le\pi$. If $L|t|\le1$, use the quadratic lower bound for $1-\cos$. If $1<L|t|\le 10$, the same bound on the subinterval of integers of size $\lfloor1/|t|\rfloor$ suffices when $|t|\le1/2$. For $|t|>1/2$ in that case, $L$ is bounded and the $j=1$ term suffices. If $L|t|>10$, the Dirichlet sum bound $|\mathbb E e^{it\Delta}|\le 1/[(2L+1)\sin(|t|/2)]$ suffices. The mixture and Taylor at low frequency now give

\begin{equation*}
1-|\phi(t)|^2\ge c V\min(t^2,L_0^{-2}),\qquad
1-|\phi(t)|^2=Vt^2(1+O(L_0^2t^2)+o(1))\quad (|t|\le c_0/L_0).
\end{equation*}

For the second claim apply Taylor on $|\Delta|\le L_0$ and use $1-\cos x\le x^2/2$ on the tail. In particular the loss at $t_*$ is $(1+o(1))V\pi^2/d^2$. Choosing $c_0$ small makes the low-frequency relative error less than $1/4$, so every diagonal phase $|t|\ge 3\pi/d$ has modulus gap at least $cV/d^2$ from principal. For $|t|\ge K\pi/d$, the lower bound on the loss in squared modulus exceeds any prescribed constant times $V/d^2$ by taking $K$ fixed sufficiently large.

\subsection{Localization relative to the principal modulus}

Write $z_j=e^{i\theta_j}$, $\bar z=d^{-1}\sum z_j$, and $D_z=\sum|z_j-\bar z|^2$. For fixed size $s$, generate a uniform subset by drawing a uniform permutation, putting $m=\min(s,d-s)$ disjoint pairs first, tossing a fair coin independently in each pair to pick one member, and including $s-m$ further predetermined positions of the permutation. By permutation symmetry this is uniform. Conditional on the permutation the modulus of its Fourier response is at most $e^{-X/8}$ where $X$ is the sum of the within-pair squared chord lengths, because each fair choice has modulus $\sqrt{1-|z_i-z_j|^2/4}$.

A uniform ordered pair of distinct indices has expected squared chord $2D_z/(d-1)$, so $\mathbb EX=2mD_z/(d-1)$. Each squared chord is at most four. For two disjoint pairs the expectation of the product of squared chords is at most $[d(d-1)/( (d-2)(d-3))]$ times the squared mean of a single chord, by allowing intersections in the numerator. Therefore $\mathbb E X^2\le C(\mathbb EX+(\mathbb EX)^2)$. The bound $1-e^{-x/8}\ge c x/(1+x)$ and Cauchy--Schwarz, in the form $\mathbb E[X/(1+X)]\ge(\mathbb EX)^2/\mathbb E[X(1+X)]$, give the fixed-size modulus loss $c\min(mD_z/d,1)$.

Since $s(d-s)/d\le m\le d$ and $\mathbb E[S(d-S)/d]=(d-1)\bar a$, averaging the fixed-size losses gives

\begin{equation*}
|R(\theta)|\le 1-c\bar a\min(D_z,1).
\end{equation*}

Indeed $\min(mD_z/d,1)\ge (m/d)\min(D_z,1)$. As $\bar a\ge V/d$, all points with $D_z\ge B V/(\bar a d^2)$ have modulus gap at least $cV/d^2$ from principal, for fixed sufficiently large $B$. At the remaining points, $D_z=O(1/d)$. Anchor at an index minimizing its distance to $\bar z$, so $\sum_j|z_j-z_a|^2\le2D_z$. Take the principal angles relative to that anchor, and subtract their mean to write $\theta=t\mathbf1+y$ with $\sum_jy_j=0$ and $\|y\|^2\le C D_z$. The sum constraint forces $dt=\pi\bmod 2\pi$. For all these points one has

\begin{equation*}
|R(t\mathbf1+y)-\phi(t)|\le \tfrac12\bar a\|y\|^2\le C B V/d^2.
\end{equation*}

This follows by centering the conditional linear term and using the exact conditional second moment $a_s\|y\|^2$. The strong radial comparison therefore excludes every $|t|\ge K\pi/d$ for fixed $K$ depending on $B$.

To handle the finitely many remaining diagonal phases, set $L=\sum_{i\in I}y_i$. If $s>d/2$, replace the subset by its complement and reverse the sign of $L$. For a uniform $m$-subset, let $p_j=(m)_j/(d)_j$ and $Q_j=\sum y_i^j$. Expansion using $Q_1=0$ gives
$\mathbb E[L^4\mid S=s]=(p_1-7p_2+12p_3-6p_4)Q_4+(3p_2-6p_3+3p_4)Q_2^2$.
Because $p_j\le p_1$ and $p_1\le C a_s$, this is at most $Ca_s\|y\|^4$. Cauchy--Schwarz against the second moment shows $\mathbb E[|L|^3\mid S=s]\le Ca_s\|y\|^3$. Moreover $s\mapsto a_s$ is $O(1/d)$-Lipschitz even when evaluated at real $s$. Thus
$\mathbb E[a_S|S-\mu|]\le a_\mu\sqrt V+CV/d\le C\bar a(\sqrt V+1)$,
using $a_\mu=\bar a+V/[d(d-1)]$ and $\bar a\ge V/d$. It follows that $a(t)/\bar a=1+O(|t|(\sqrt V+1))$ for $|t|\le K\pi/d$, by centering phase at $\mu$ and $|e^{ix}-1|\le |x|$. In this range $|\phi(t)|=1-o(1)$ and the logarithmic expansion is

\begin{equation*}
\log[R(t\mathbf1+y)/\phi(t)]=-\tfrac12 a(t)\|y\|^2+O(\bar a\|y\|^3).
\end{equation*}

The bound holds uniformly for $\|y\|=O(d^{-1/2})$, by Taylor and $\bar a\le1$. Consequently the local real part is at most $-c\bar a\|y\|^2$. Together the estimates give a relative modulus gap at least $c\min(V/d^2,\bar a/d^2)$ outside the principal balls $\|y\|\le d^{-1}$. This is at least $c d^{-5}$ in the first regime and $c_\eta d^{-3}$ in the second, by the variance bounds and $\bar a\ge1/(11d)$.

\subsection{Integrating the two principal balls}

The translated torus has period lattice $2\pi\{m\in\mathbb Z^d:\sum_jm_j=0\}$ in $\mathbf1^\perp$. The basis $2\pi(e_j-e_d)$ has Gram determinant $d$ after removal of the $2\pi$ factors, so the Euclidean covolume is $\sqrt d(2\pi)^{d-1}$. The two balls are disjoint and inject into the torus for large $d$. Normalizing the full Euclidean complex Gaussian gives one contribution $\phi(t_*)^n(2\pi n)^{-(d-1)/2}a(t_*)^{-(d-1)/2}/\sqrt d$ and its conjugate. The correction in the logarithm is $O(n\bar a \|y\|^3)$. On the ball this is absorbed by weakening the real quadratic by a relative $O(d^{-1})$, which changes the Gaussian normalization by a bounded factor. Also $(|a|/\mathrm{Re}\,a)^{(d-1)/2}=\exp(O(1))$ since $a/\bar a=1+O(d^{-1/2})$. The Gaussian third moment shows the relative integration error is $O(d^{3/2}/\sqrt{n\bar a})=o(1)$, via $|e^z-1|\le |z|e^{(\mathrm{Re}\,z)_+}$. The discarded Gaussian mass outside $d^{-1}$ is exponentially small since $n\bar a d^{-2}\gg d$. The local error is $O(d^2/\sqrt n)$ by $\bar a\ge1/(11d)$. In the first regime the relative global tail is at most $\exp(C_1d\log n-c_1 n d^{-5})$. In the second it is at most $\exp(C_1d\log n-c_\eta n d^{-3})$. For either regime, write $\ell=5$ or $3$, respectively. At $n=C d^{\ell+1}\log(d+1)$ the ratio $n/(d^{\ell+1}\log n)$ tends to $C/(\ell+1)$. Since $x/\log x$ increases for $x>e$, a sufficiently large fixed $C$ makes the global exponent tend to $-\infty$ uniformly throughout the corresponding region. This completes the absolute error estimate in both regimes.

\section{Upper-edge two-saddle approximation with growing cap}\label{sec:growing-cap-upper-saddle}

2026-11-26 \enspace Private research note

There are absolute constants $h_0>0$ and $C_*>0$ with the following property. Let $d\to\infty$, $1\le r<d$, and $n\ge C_*d^4\log(d+1)$. Take $t>0$ and let $S$ have probabilities proportional to $\binom ds t^s$ for $0\le s\le r$. With $J=r-S$, assume

\begin{equation*}
h=\mathbb E J\le h_0,\qquad nh/d^2\longrightarrow\infty.
\end{equation*}

Let $H(w)=\sum_{s\le r}\binom ds w^s$, $\phi=\pi/d$, $H_\tau=H(te^{i\tau})$, and

\begin{equation*}
\lambda_s=\frac{s(d-s)}{d(d-1)},\qquad
\Lambda_\tau=\frac{\mathbb E(\lambda_S e^{-i\tau J})}{\mathbb E e^{-i\tau J}},
\qquad \Lambda_0=\mathbb E\lambda_S.
\end{equation*}

For the balanced-array polynomial $P_n$, the uniform approximation in this regime is

\begin{equation}
P_n(-t^d)=2\mathrm{Re}\left(
 \frac{H_\phi^n \Lambda_\phi^{-(d-1)/2}}{\sqrt d(2\pi n)^{(d-1)/2}}
 \right)
 +o\left(\frac{|H_\phi|^n|\Lambda_\phi|^{-(d-1)/2}}{\sqrt d(2\pi n)^{(d-1)/2}}\right).
\tag{S1}
\end{equation}

The power uses the logarithm in the right half-plane. The error is absolute relative to the magnitude of one saddle, including when the main real part is small. The condition $nh/d^2\to\infty$ may hold arbitrarily slowly.

\subsection{Tilt and diagonal gaps}

The ratio of successive unnormalized deficit weights at $j,j+1$ is $(r-j)/[(d-r+j+1)t]$, decreasing in $j$. The first ratio $q$ satisfies $q=\Pr(J=1)/\Pr(J=0)\le h/(1-h)$. Thus $\Pr(J=j)\le q^j$ and $\mathbb E J^\ell=O_\ell(h)$ for every fixed positive integer $\ell$, uniformly once $h_0$ is sufficiently small. Also $h\le Cq$ by the same bound, so $\Pr(J=1)\asymp h$ and $\mathrm{Var}(J)\asymp h$. For $j\le r$,

\begin{equation*}
\frac{\lambda_{r-j}}{\lambda_r}
=\frac{r-j}{r}\frac{d-r+j}{d-r}\le1+j.
\end{equation*}

As $\Pr(J=0)\ge1-h$, this implies $\Lambda_0\asymp\lambda_r$, $c/d\le\Lambda_0\le C$, and $\mathbb E[\lambda_S^a J^\ell]=O_{a,\ell}(\Lambda_0^a h)$ for fixed $a,\ell\ge1$. The same bound holds for the expectation of $\lambda_S^a\mathbf1_{J>0}$.

Write $g(\tau)=\mathbb E e^{-i\tau J}$, so $H_\tau/H(t)=e^{ir\tau}g(\tau)$. Uniformly on $[-\pi,\pi]$, $g(\tau)=1+O(h)$ and $|g(\tau)|\le\exp[-ch(1-\cos\tau)]$, the latter by retaining the $J=0,1$ terms in the independent-copy formula for $1-|g|^2$. Centering $J$ and differentiating its characteristic function gives $(\log|g(\tau)|)''=-\mathrm{Var}(J)+O(h|\tau|)$ near zero by the centered third-moment bound. Evenness gives $\log|g(\tau)/g(\phi)|\le-ch(\tau^2-\phi^2)$ for $\phi\le|\tau|\le\eta_0$ with fixed sufficiently small $\eta_0$, and $|g(\phi)|\ge e^{-Ch/d^2}$. Beyond $\eta_0$, the global diagonal gap and the lower bound at $\phi$ give the same comparison with another constant. Thus

\begin{equation}
\log|H_\tau/H_\phi|\le-ch(\tau^2-\phi^2)
\qquad(\phi\le|\tau|\le\pi).
\tag{S2}
\end{equation}

For later volume comparisons, differentiating the numerator and denominator of $\Lambda_\tau$ once or twice costs $O(\Lambda_0h)$ and $O(h)$, respectively. This uses the $J$-moment bounds and $|g(\tau)|\ge1-O(h)$. Uniformly, $\Lambda_\tau=\Lambda_0(1+O(h))$, as follows by separating the $J=0$ term in numerator and denominator. Conjugation makes the first derivative at zero purely imaginary. Hence

\begin{equation}
\mathrm{Re}\Lambda_\tau/\Lambda_0=1+O(h\min(\tau^2,1)),\qquad
|\mathrm{Im}\Lambda_\tau|/\Lambda_0=O(h\min(|\tau|,1)).
\tag{S3}
\end{equation}

In particular $|\Lambda_\phi|/\Lambda_0=1+O(h/d^2)$.

\subsection{Localization relative to the principal modulus}

Let $R(\mathbf z)=\sum_{|A|\le r}t^{|A|}\prod_{i\in A}z_i$ on the torus with $\prod_i z_i=-1$, with normalized Haar measure. The integral of $R^n$ is $P_n(-t^d)$ by character orthogonality. Put $\overline z=d^{-1}\sum_i z_i$ and $D_z=\sum_i|z_i-\overline z|^2$.

For fixed size $s$ let $m=\min(s,d-s)$. Generate a uniform subset by putting $m$ disjoint random pairs in the first $2m$ positions of a uniform permutation, choosing one element from each pair by independent fair coins, and including the next $s-m$ positions. By symmetry the subset is uniform. Conditional on the permutation the modulus of its Fourier response is at most $e^{-X/8}$ where $X$ is the sum of the squared chords within the pairs. Indeed the squared modulus of the mean of the two points of a pair is one minus one fourth of the squared chord. For one ordered random pair the mean squared chord is $2D_z/(d-1)$, so $\mathbb E X=2mD_z/(d-1)$. Each squared chord is at most four. For two disjoint pairs their product expectation is at most $d(d-1)/[(d-2)(d-3)]$ times the squared one-pair expectation, by discarding the disjointness restriction in the numerator before dividing by the number of disjoint pairs. Thus $\mathbb E X^2\le C(\mathbb EX+(\mathbb EX)^2)$. From Cauchy--Schwarz,

\begin{equation*}
\mathbb E\frac{X}{1+X}\ge
\frac{(\mathbb E X)^2}{\mathbb E(X+X^2)}.
\end{equation*}

Since $1-e^{-x/8}\ge c x/(1+x)$, the fixed-size response has modulus at most $1-c\min(mD_z/d,1)$. This also holds for $m=0$. Averaging over $S$ and using $m\ge s(d-s)/d$ gives

\begin{equation}
|R/H(t)|\le1-c\Lambda_0\min(D_z,1).
\tag{S4}
\end{equation}

Here $\min(mD_z/d,1)\ge(m/d)\min(D_z,1)$ and $\mathbb E[S(d-S)/d]=(d-1)\Lambda_0$.

For $D_z\ge1/d$, (S4) gives modulus at most $H(t)(1-c/d^2)$. By reducing $h_0$ this is at most $|H_\phi|(1-c'/d^2)$, since $|H_\phi|/H(t)\ge e^{-Ch/d^2}$. Its absolute integral, divided by the principal scale

\begin{equation*}
|H_\phi|^n I_0,\qquad I_0=(2\pi n\Lambda_0)^{-(d-1)/2}/\sqrt d,
\end{equation*}

is bounded by $\exp(-c'n/d^2+O(d\log n))$, which is negligible in the asserted range.

For $D_z<1/d$, choose a pin minimizing $\sum_i|z_i-z_j|^2$, which is at most $2D_z$ by averaging. Lift the angles relative to this pin, take their average $\tau$ modulo $2\pi$ in $(-\pi,\pi]$, and write $z_i=e^{i(\tau+y_i)}$ with $\sum_i y_i=0$ and $\|y\|^2=O(1/d)$. The norm bound follows from the small-chord bound and subtracting the mean of the lifted differences. The common phase is an odd multiple of $\phi$ modulo $2\pi$. The representation is unique within this small norm domain: if two lifts give the same point, their relative shifts by integer multiples of $2\pi$ agree at all coordinates because the pairwise angular differences are small, and averaging fixes the canonical phase. In free coordinates the measure is $(2\pi)^{-(d-1)}dy_1\cdots dy_{d-1}$, with $y_d=-\sum_{i<d}y_i$. Each such chart contains the full ball $\|y\|\le d^{-1}$, because $D_z\le\sum_i|z_i-e^{i\tau}|^2\le \|y\|^2$.

\subsection{Transverse expansion at all common phases}

For $L=\sum_{i\in A}y_i$ with $A$ a uniform size-$s$ subset, write $P_j=\sum_i y_i^j$ and $m_j=(s)_j/(d)_j$. Conditional means vanish. Expansion using $P_1=0$ gives

\begin{equation*}
\begin{aligned}
\mathbb E(L^2\mid S=s)&=\lambda_s P_2,\\
\mathbb E(L^3\mid S=s)&=\lambda_s\frac{d-2s}{d-2}P_3,\\
\mathbb E(L^4\mid S=s)&=(m_1-7m_2+12m_3-6m_4)P_4
                  +3(m_2-2m_3+m_4)P_2^2.
\end{aligned}
\end{equation*}

For the cubic, the mixed sum $\sum_{i\ne j}y_i^2y_j$ is $-P_3$ and $e_3(\mathbf y)=P_3/3$. For the quartic the sums with respective patterns $3+1$, $2+2$, $2+1+1$, and $1+1+1+1$, with equal exponents unlabeled, are $-P_4$, $(P_2^2-P_4)/2$, $P_4-P_2^2/2$, and $P_2^2/8-P_4/4$. The coefficient of $P_4$ is $O(\lambda_s)$ by writing it in terms of $m_j-m_{j+1}=m_j(d-s)/(d-j)=O(\lambda_s)$. The coefficient of $P_2^2$ is $O(\lambda_s^2)$ by $m_2-2m_3+m_4=(s)_2(d-s)_2/(d)_4$. In particular $\mathbb E(L^4\mid S=s)\le C\lambda_s\|y\|^4$. Cauchy--Schwarz against the conditional second moment gives $\mathbb E(|L|^3\mid S=s)\le C\lambda_s\|y\|^3$.

Normalize the complex size weights by $g(\tau)$ and tilt by $e^{-i\tau J}$. Its absolute weights are bounded by a constant times the underlying real probabilities, uniformly on the entire circle. The conditional linear mean vanishes. Taylor's bound for the real exponential through degree two gives

\begin{equation}
\log(R/H_\tau)=-\tfrac12\Lambda_\tau\|y\|^2+O(\Lambda_0\|y\|^3)
\tag{S5}
\end{equation}

throughout the localized region. In detail, the exponential remainder is bounded by the conditional third absolute moment, and $|R/H_\tau-1|=O(\Lambda_0\|y\|^2)$ there by the vanishing linear term. The error from taking the logarithm is thus absorbed in (S5), with the branch near one. By (S3), $\mathrm{Re}\Lambda_\tau\ge(1-Ch_0)\Lambda_0$ for all phases. For small fixed $h_0$, (S5) gives real transverse decay at least $c\Lambda_0\|y\|^2$. Thus, in the retained region but outside the balls of radius $d^{-1}$ about all common phases, $|R|\le|H_\phi| e^{-c/d^3}$ by (S2). The absolute integral there relative to $|H_\phi|^n I_0$ is at most $\exp[-cn/d^3+O(d\log n)]$. At $n=C_* d^4\log(d+1)$, the ratio $n/(d^4\log n)$ tends to $C_*/4$. This ratio increases with $n$ in the stated region. Taking $C_*$ sufficiently large makes the last error negligible uniformly. The same choice also covers the earlier global tail.

\subsection{Summing the phase balls and evaluating the principal integrals}

Inside a ball $\|y\|\le d^{-1}$, (S3) and (S5) bound the absolute integrand by a Gaussian with quadratic coefficient $\Lambda_0[1-O(d^{-1}+h\min(\tau^2,1))]$. Its integral relative to $I_0$ is at most $\exp[O(1+dh\min(\tau^2,1))]$. The zero-sum Gaussian determinant in the free coordinates is $d$, since the quadratic form matrix is $I+\mathbf1\mathbf1^\top$. For every nonprincipal phase $|\tau|\ge3\phi$ the factor involving $dh\tau^2$ is absorbed by half the gap in (S2), because $\tau^2\le(9/8)(\tau^2-\phi^2)$ and $n\gg d$. Summing the odd phase labels gives relative error at most

\begin{equation*}
C\sum_{j\ge1}\exp[-c(nh/d^2)((2j+1)^2-1)]
=O(e^{-c'nh/d^2})=o(1).
\end{equation*}

There is no prefactor growing with $d$. This is the reason the condition $nh/d^2\to\infty$ can be arbitrarily slow.

At $\tau=\phi$, compare $\exp(n\log(R/H_\phi))$ with $\exp(-n\Lambda_\phi\|y\|^2/2)$ by $|e^A-e^B|\le |A-B|e^{\max(\mathrm{Re} A,\mathrm{Re} B)}$. The weakened Gaussian just used has normalization $O(I_0)$ at $\phi$. Its normalized third norm moment is $O(d^{3/2}/(n\Lambda_0)^{3/2})$. Consequently the integrated absolute error of (S5) is at most

\begin{equation*}
O(I_0\,d^{3/2}/\sqrt{n\Lambda_0})
=O(I_0\,d^2/\sqrt n)
=o(I_0).
\end{equation*}

The comparison Gaussian extends to the full hyperplane at cost $o(I_0)$, since its absolute tail outside $d^{-1}$ is bounded relative to $I_0$ by $C2^{d/2}\exp(-c n\Lambda_0/d^2)$ using the Gaussian quadratic exponential moment. Its complex integral is $(2\pi n\Lambda_\phi)^{-(d-1)/2}/\sqrt d$ and its modulus is $(1+o(1))I_0$ by (S3). The ball at $-\phi$ gives the conjugate integral. This completes (S1).

\end{document}
